\section{Discussion}

We have shown that the long-time failure of force-free molecular dynamics propagators is not
primarily a one-step accuracy problem. A published propagator with a fifteen-fold larger pointwise
error reproduces the canonical water ensemble, while a pointwise-better model develops a large mode
imbalance and fails the same NVT screen. The failure is also not an energy-drift problem: the model
that heats most strongly in NVE passes NVT, and the most NVE-stable model fails. The discriminating
object is the mode-resolved stationary measure of the composed propagator plus thermostat. This
diagnostic is, to our knowledge, the first systematic demonstration of the pointwise--ensemble
decoupling for force-free MD, and it changes what a training or acceptance criterion should
measure.

The remedy follows from the diagnosis. A frozen backbone supplies the structure---the displacement
that determines the position marginal---while a mode-resolved generalized-Langevin readout supplies
the velocity that enters the next backbone call. The readout fits a per-mode AR(\(p\)) memory
kernel to the reference velocity autocorrelation, injects noise with the fluctuation--dissipation
consistent variance, and rescales each mode or group to its canonical kinetic-energy target at every step.
The partition is guaranteed by construction; the memory kernel matches the multi-lag dynamics that
a single gain cannot. The backbone is not retrained, and the readout is fitted from a short
reference NVT trajectory. Importantly, we do not claim that the force-free backbone alone
reproduces long-time dynamics; the raw backbones fail the NVT screen (Table~3). The claim is that
the fixed readout corrects the stationary measure of a frozen backbone. On water, the result is a model that is stable on two independent initial
conditions, initial-condition independent in its partition, ACF, RDF, and VDOS, and more accurate
than the published checkpoint on the RDFs and VDOS in our analysis.

% Table 2 — Conceptual comparison with the closest prior art.
\begin{table}[t]
\centering
\small
\caption{What the GLR controls compared with the closest prior art. The distinction is the
object (a readout on a frozen force-free propagator) and the target (the canonical mean mode/group
kinetic partition, with a stochastic marginal option, plus the multi-lag velocity ACF), not the
use of a memory kernel per se.}
\label{tab:prior}
\resizebox{\linewidth}{!}{%
\begin{tabular}{@{}p{2.1cm}p{3.1cm}p{4.0cm}p{3.6cm}@{}}
\toprule
Approach & Object & What is controlled or learned & Evidence \\
\midrule
TrajCast & force-free propagator + global CSVR & one-step displacement/velocity; total
temperature & water RDF/VDOS/diffusion \\
Bigi--Ceriotti & deterministic structure-preserving map & symplectic/time-reversible structure;
energy/equipartition & water, GeTe, Ar, Al \\
MAMD & LSTM velocity update & stability via periodic learned velocity injection & several MD
systems \\
GLE thermostat & memory kernel + colored noise in an integrator & FD stationary distribution of
the integrator & generic MD \\
\textbf{GLR} & \textbf{readout on a frozen force-free propagator} & \textbf{mean mode kinetic
partition (+ stochastic marginal) + multi-lag velocity ACF} & \textbf{water NVT, two 2.5 ns
starts; quartz 5 fs, two seeds} \\
\bottomrule
\end{tabular}}
\end{table}


The relationship to prior work is complementary rather than competitive. TrajCast established that
force-free propagators can reproduce water physics \cite{trajcast}; our work explains when that
success is fragile and adds a readout that makes the mode partition a design constraint rather
than a lucky outcome. Structure-preserving action maps enforce symplectic or time-reversible
structure and improve energy conservation and equipartition \cite{bigi_action}; they target
mechanical structure, whereas the GLR targets the canonical NVT measure and the multi-lag velocity
statistics. The two could be combined: an action map would supply a better backbone, and the GLR
would enforce the mode partition on top of it. Learned velocity updates recover stability with a
black-box LSTM and periodic injection \cite{mamd}; the GLR is interpretable, mode-resolved, and
tied to the reference ACF, and it guarantees the partition rather than only delaying divergence.
GLE thermostats provide the same memory-kernel and colored-noise physics
\cite{ceriotti_prl,ceriotti_jctc}, and frequency-resolved thermostats can control selected modes
\cite{zpe_thermostat}; the difference is the object. A GLE thermostat acts on the forces and
velocities of an MD integrator. The GLR is a readout on a frozen force-free propagator, fitted from
the reference ACF, with an explicit group-resolved variance guarantee and a diagnostic that motivates it.
Stochastic and generative propagators offer a different route to equilibrium
\cite{timewarp,mdgen}, at the cost of MCMC, an energy function, or a larger generative model; the
GLR is a lightweight post-hoc readout.

The diagnostic has implications beyond this method. Force-free propagators are usually selected and
reported by one-step MAE, sometimes with a short rollout score. Our results show that this
criterion can rank models in the wrong order. An ensemble-level, mode-resolved
evaluation---kinetic partition, RDF, VDOS, and diffusion---is necessary to detect the failure. The
same lesson applies to other learned dynamical maps: a low conditional-mean error does not imply a
correct invariant measure, and a stable rollout does not imply a correct one. We therefore
recommend that force-free MD papers report the mode partition and at least one long-time structural
or dynamical observable, not only one-step errors.

Several limitations define the boundary of the claim. First, the analytic guarantee covers the
\emph{mean} mode kinetic partition and the velocity ACF, not the full joint position--velocity
distribution. The deterministic per-step rescale fixes the mode kinetic energy to its canonical
mean and suppresses its equilibrium fluctuations (total KE std 0.010 eV versus 0.440 eV in the
reference). A stochastic-target variant restores the KE variance (0.437/0.429 eV), but its
Green--Kubo integral (\(1.16\times10^{-9}\)) is inconsistent with the MSD (\(1.44\times10^{-9}\)):
the extra fluctuations modify the negative translation-ACF tail. The deterministic target is
therefore used for production because it preserves GK--MSD consistency; the joint
KE-variance/transport problem remains open. The marginal velocity distribution is nevertheless
close to the reference (component kurtosis 2.93/5.49/5.93 versus 3.04/5.62/5.57). The position
marginal is learned by the backbone and validated empirically; the stochastic rescale guarantees
the kinetic-energy marginal, not the full joint position--velocity distribution, which would
require an energy-based sampler such as MCMC. Second,
the readout requires a short reference NVT trajectory to fit the ACF. This is much cheaper than
retraining and is available for most systems of interest, but it is a data requirement, not a free
lunch. The fit uses only the mode velocity ACF; the RDF, VDOS, and diffusion are held out and are
not used to fit any readout parameter. Relatedly, the FD-noise ablation is negative at
\(\lambda=0.1\): removing the noise does not change the partition, RDF, VDOS, or \(D\). The memory
kernel and the rescale are the load-bearing components; the noise is retained for theoretical
consistency, not because the data require it. Third, we demonstrate one water model, two backbone
architectures (relational and phase-space), and a second material (quartz). The readout is
empirically backbone-agnostic across the two water backbones and transfers to a covalent solid
through the variance-rescale component; the full memory kernel is water-specific because the 5 fs
quartz backbone already has the correct mode ACF. A broader material sweep is still needed to
establish cross-system generality. Fourth, the internal ACF at 500 fs remains too persistent for AR(10); AR(20) reduces
the discrepancy but does not eliminate it. Fifth, the diffusion coefficient is 26\% below the
reference on both 2.5 ns starts (1.39--1.41 versus 1.88\(\times10^{-9}\)). The Green--Kubo integral
shows that the deficit is in the translation ACF, and the decoupled position update prevents a
simple blend or displacement correction from closing it; a position-consistent readout or an
end-to-end centre-of-mass head is future work. Sixth, the
backbone uses a local 6 \AA{} cutoff and no
long-range electrostatics; the GLR does not address long-range physics, and the water results
should not be read as a claim about large-scale dielectric or interfacial behaviour. Seventh, a
5 fs phonon-band extension to quartz is stable. Without the readout the 5 fs backbone cools to
235 K; the stochastic phonon-band variance rescale (two seeds) sets 299.9 K, reproduces the total
kinetic-energy variance (0.397/0.399 versus 0.4017 eV), improves the Si--O RDF from 0.271 to
0.028, and preserves the phonon VDOS at 0.950/0.955. A coarser species partition gives 301.68 K, exact
Si/O equipartition, total KE std 0.4012 eV, and VDOS 0.898/0.955. The deterministic rescale is an
ablation: it fixes the mean but suppresses the kinetic-energy variance. The full species GLR
(memory + noise) is a negative control: species velocity is not a phonon eigenbasis, and the
readout collapses the VDOS to 0.07 even at 5 fs. A 30 fs backbone shows the same pattern more
strongly. A 486-mode phonon-basis GLR is also negative: the basis is canonical, but the 5 fs
backbone's mode ACF is already correct, so the AR innovations broaden the phonon lines (VDOS
0.734/0.838). The variance-rescale component therefore generalizes across a covalent solid; the
memory kernel is useful only when the backbone's mode ACF is wrong, and its basis must be
physically meaningful. Finally, the 2.5 ns
production runs use one seed per start; the 100 ps seed study is reported in the Results, and a
second 2.5 ns seed would tighten the transport error.

The outlook is straightforward. The readout can be trained jointly with the backbone, rather than
fitted post hoc, which may improve the position marginal while retaining the partition guarantee.
The memory kernel can be made system- and temperature-adaptive, and the mode decomposition can be
extended to phonons or other collective coordinates. Combined with a structure-preserving backbone
and a long-range operator, the same diagnostic-and-readout strategy could make force-free
propagators reliable beyond water.
